% Generated draft: requires mathematical and source review.
\documentclass[11pt]{article}
\usepackage[margin=25mm]{geometry}
\usepackage{fontspec}
\setmainfont{Latin Modern Roman}
\usepackage{amsmath,amssymb,mathtools,unicode-math}
\setmathfont{Latin Modern Math}
\usepackage{enumitem,needspace,xurl,hyperref}
\setlength{\parindent}{0pt}
\setlength{\parskip}{5pt}
\setlength{\emergencystretch}{3em}
\newcommand{\N}{\mathbb{N}}
\newcommand{\Z}{\mathbb{Z}}
\newcommand{\Q}{\mathbb{Q}}
\newcommand{\R}{\mathbb{R}}
\newcommand{\C}{\mathbb{C}}
\newcommand{\F}{\mathbb{F}}
\newcommand{\A}{\mathbb{A}}
\newcommand{\PP}{\mathbb{P}}
\newcommand{\E}{\mathbb{E}}
\newcommand{\eps}{\varepsilon}
\newcommand{\dd}{\,\mathrm{d}}
\begin{document}
\section*{UnsolvedMath ID: 30004563}
\textbf{Source problem code:} catalogue entry
\subsection{Definitions and mathematical statement}
The supplied record asks about \(\alpha\)-immersions with \(\alpha>1\) and the potential central points arising from a cube move. It does not define an \(\alpha\)-immersion, the admissible local data for a cube move, or the conditions making a central point potential. Consequently, a complete geometric formulation cannot be reconstructed from this record alone.

For a conditional formulation, write \(\mathcal D_\alpha\) for the class of admissible cube-move inputs under the conventions of the cited sources, and write \(C_\alpha(D)\) for the set of potential central points associated with an input \(D\in\mathcal D_\alpha\). These are placeholders for definitions requiring source verification, not additional geometric assumptions. Interpreting the count as the number of distinct potential central points, the proposed universal upper bound is
\[
\forall\,\alpha>1\;\forall\,D\in\mathcal D_\alpha,
\qquad \#C_\alpha(D)\leq 3.
\]
Here the parameter domain and all admissibility, boundary, and degeneracy conditions must be supplied from the cited construction; they are not specified in the record.

The question is whether three is the maximal possible number. Sharpness across the entire parameter range would additionally mean
\[
\exists\,\alpha>1\;\exists\,D\in\mathcal D_\alpha,
\qquad \#C_\alpha(D)=3.
\]
Sharpness for every fixed parameter would instead mean
\[
\forall\,\alpha>1\;\exists\,D\in\mathcal D_\alpha,
\qquad \#C_\alpha(D)=3.
\]
The supplied wording does not distinguish these sharpness quantifiers. It also does not specify whether coincident candidates are counted separately, whether candidates must give admissible resulting immersions, or which degenerate configurations are permitted. These qualifications remain unresolved.

The supplied dated assessment reports existence for \(\alpha>1\), configurations with up to three potential central points, and a conjectural upper bound of three without a general proof. These are supplied literature claims, not independently verified mathematical conclusions.
\subsection{Short English statement}
For admissible cube moves of \(\alpha\)-immersions with \(\alpha>1\), must there be at most three potential central points, and is this bound attained? The geometric definitions, counting conventions, and whether attainment is required for each fixed \(\alpha\) are missing from the supplied record. This is a conditional statement-only draft requiring mathematical review.
\subsection{Sources}
\begin{itemize}
\item Supplied local UnsolvedMath registry record, permanent ID \(30004563\), title ``Maximum Central Points in Cube-Move \(\alpha\)-Immersions,'' exact source problem code \texttt{catalogue entry}, in the Oberwolfach Reports Open Problems collection. No individual proposer is supplied; the proposed year is recorded as 2020.
\item Supplied report citation: \emph{Mini-Workshop: Dimers, Ising and Spanning Trees beyond the Critical Isoradial Case} (2021), Oberwolfach Reports, pp.\ 1759--1789, DOI \texttt{10.4171/owr/2020/35}. \url{https://doi.org/10.4171/owr/2020/35}. The record describes Oberwolfach Reports as documenting workshops and their open-problem discussions.
\item Supplied literature source: \emph{Cube moves for s-embeddings and alpha-realizations}. \url{https://arxiv.org/abs/2003.08941}. Authors and a specific theorem, conjecture, or page locator are not supplied.
\item Supplied literature source: ``OWR 2020/35 report.'' \url{https://doi.org/10.4171/owr/2020/35}.
\item Literature assessment date: \texttt{2026-08-22}. The supplied dated curation labels the problem open and reports existence for \(\alpha>1\), configurations with up to three potential central points, and no general proof that three is maximal. The background explicitly calls this dated literature triage derived from the supplied OWR curation and requests independent verification of time-sensitive status and cited claims. This assessment has not been independently verified here.
\item The registry labels the statement \texttt{corrected\_verified} and reports that it was checked against the official Oberwolfach report, preserving its hypotheses, notation, and open target. This is supplied verification metadata, not a claim of source consultation or verification in this draft.
\end{itemize}
\end{document}
